% !TEX root = ../../../main.tex
% 所属：第一部分：单变量分析学 / 第12章：函数项序列与函数项级数
% 节：12.5 Taylor级数展开
% 正文按原文件顺序保存；由 ../chapter12.tex 调用。
\section{Taylor级数展开}

这一节讨论一个函数何时可以展开成以 $x_0$ 为中心的幂级数.设在 $x_0$ 的某个邻域内有
\[
\left|x-x_0\right|<\rho,
\]
并且
\[
f(x)=a_0+a_1(x-x_0)+\cdots+a_n(x-x_0)^n+\cdots .\tag{*}
\]

$f(x)$能够展开成(*)形式的必要条件：
\begin{enumerate}
	\item $f\in C^\infty(x_0-\rho,x_0+\rho)$，且系数只能取
	\[
	a_n=\frac{f^{(n)}(x_0)}{n!};
	\]
	注意, $f\in C^\infty(x_0-\rho,x_0+\rho)$只是能展开成幂级数的必要条件,不能保证函数一定能展开成幂级数.
	
	由$a_n=\dfrac{f^{(n)}(x_0)}{n!}$我们成这一种展开为Taylor级数
	\item 所得 Taylor 级数收敛;
	\item 其和函数确为 $f(x)$.
\end{enumerate}
\begin{example}
	设
	\[
	f(x)=
	\begin{cases}
		e^{-\frac{1}{x^2}},&x\neq0,\\[6pt]
		0,&x=0.
	\end{cases}
	\]
\end{example}

\begin{solution}
	当 $x\neq0$ 时,
	\[
	f'(x)=\frac{2}{x^3}e^{-\frac{1}{x^2}},
	\]
	\[
	f''(x)=\left(\frac{4}{x^6}-\frac{6}{x^4}\right)e^{-\frac{1}{x^2}},
	\]
	\[
	\cdots\cdots
	\]
	\[
	f^{(n)}(x)=P_{3n}\left(\frac{1}{x}\right)e^{-\frac{1}{x^2}},
	\]
	\[
	\cdots\cdots
	\]
	其中 $P_m(u)$ 是关于 $u$ 的 $m$ 次多项式.事实上,若
	\[
	f^{(n)}(x)=P_{3n}\left(\frac{1}{x}\right)e^{-\frac{1}{x^2}},
	\]
	则求导后仍可写为
	\[
	f^{(n+1)}(x)=P_{3n+3}\left(\frac{1}{x}\right)e^{-\frac{1}{x^2}},
	\]
	于是上式由归纳得到.
	
	由此可以依次得到
	\[
	\begin{aligned}
		f'(0)
		&=\lim_{x\to0}\frac{f(x)-f(0)}{x}
		=\lim_{x\to0}\frac{1}{x}e^{-\frac{1}{x^2}}
		=0,\\[6pt]
		f''(0)
		&=\lim_{x\to0}\frac{f'(x)-f'(0)}{x}
		=\lim_{x\to0}\frac{2}{x^4}e^{-\frac{1}{x^2}}
		=0,
	\end{aligned}
	\]
	\[
	\cdots\cdots
	\]
	并且一般地,
	\[
	\begin{aligned}
		f^{(n)}(0)
		&=\lim_{x\to0}\frac{f^{(n-1)}(x)-f^{(n-1)}(0)}{x}\\
		&=\lim_{x\to0}P_{3n-2}\left(\frac{1}{x}\right)e^{-\frac{1}{x^2}}
		=0.
	\end{aligned}
	\]
	这里用到的是:任意多项式增长都被 $e^{\dfrac{1}{x^2}}$ 压过,故
	\[
	\lim_{x\to0}P\left(\frac{1}{x}\right)e^{-\frac{1}{x^2}}=0.
	\]
	
	因此 $f(x)$ 在 $x=0$ 的 Taylor 级数为
	\[
	0+0x+\frac{0}{2!}x^2+\frac{0}{3!}x^3+\cdots+\frac{0}{n!}x^n+\cdots,
	\]
	它在 $(-\infty,+\infty)$ 上收敛于和函数
	\[
	S(x)=0.
	\]
	显然,当 $x\neq0$ 时,
	\[
	S(x)\neq f(x).
	\]
	这说明,一个任意阶可导的函数的 Taylor 级数并非一定能收敛于函数本身.
\end{solution}

\begin{definition}
	若一个函数在某邻域内等于其 Taylor 级数且满足上述的1,2,3条,则称其在该邻域内为解析函数,记为 $f\in C^\omega$.
\end{definition}


由上面的必要条件,若 $f$ 可以展开为幂级数,则其部分和函数必为
\[
S_n(x)=f(x_0)+f'(x_0)(x-x_0)+\cdots+\frac{f^{(n)}(x_0)}{n!}(x-x_0)^n.
\]
这称为 Taylor 展开式的不带余项形式.于是
\[
f(x)=S_n(x)+r_n(x),
\]
其中 $r_n(x)$ 称为余项.若在 $\left|x-x_0\right|<\rho$ 内有
\[
r_n(x)\to0,
\]
则 $f(x)$ 等于其 Taylor 级数.
\begin{proposition}
	(Taylor 公式的积分余项)
	若 $f\in C^{n+1}(x_0-\rho,x_0+\rho)$,则
	\[
	f(x)=\sum_{k=0}^{n}\frac{f^{(k)}(x_0)}{k!}(x-x_0)^k+r_n(x),
	\]
	其中
	\[
	r_n(x)=\frac{1}{n!}\int_{x_0}^{x}f^{(n+1)}(t)(x-t)^n\dd t.
	\]
\end{proposition}

\begin{proof}
	令
	\[
	r_n(x)=f(x)-S_n(x).
	\]
	则
	\[
	r_n(x_0)=r_n'(x_0)=\cdots=r_n^{(n)}(x_0)=0,
	\]
	且
	\[
	r_n^{(n+1)}(x)=f^{(n+1)}(x).
	\]
	由微积分基本公式及分部积分,
	\[
	\begin{aligned}
		r_n(x)
		&=r_n(x_0)+\int_{x_0}^{x}r_n'(t)\dd t = \int_{x_0}^{x}r_n'(t)\dd t   =  \int_{x_0}^{x}r_n'(t)\dd (t-x)  \\
		&=r'_n(x_0)(t-x)|_{x_0}^{x} -\int_{x_0}^{x}r_n''(t)(t-x)\dd t=\int_{x_0}^{x}r_n''(t)(x-t)\dd t                                  \\
		&=\frac{1}{2!}\int_{x_0}^{x}r_n''(t)\dd (x-t)^2  =\frac{1}{2!}\int_{x_0}^{x}r_n^{(3)}(t)(x-t)^2\dd t                 \\
		&=\cdots \\
		&=\frac{1}{n!}\int_{x_0}^{x}r_n^{(n+1)}(t)(x-t)^n\dd t=\frac{1}{n!}\int_{x_0}^{x}f^{(n+1)}(t)(x-t)^n\dd t.
	\end{aligned}
	\]
\end{proof}

(Lagrange余项)
	在上面的条件下,存在 $\xi$ 介于 $x_0$ 与 $x$ 之间,使得
	\[
	r_n(x)=\frac{f^{(n+1)}(\xi)}{(n+1)!}(x-x_0)^{n+1}.
	\]


(Cauchy余项)
	在上面的条件下,存在 $\xi$ 介于 $x_0$ 与 $x$ 之间,使得
	\[
	r_n(x)=\frac{f^{(n+1)}(\xi)}{n!}(x-\xi)^n(x-x_0).
	\]


下面总结 Taylor 级数展开的常用方法.
\begin{enumerate}
	\item 直接计算
	\[
	a_n=\frac{f^{(n)}(x_0)}{n!}
	\]
	以及余项 $r_n(x)$,证明 $r_n(x)\to0$,再代入 Taylor 级数,最后判断收敛域.
	\item 由已知展开式出发,通过代换、加减乘除等运算,可以得到新的 Taylor 展开式.
	\item 若已知幂级数在相应区间内可以逐项积分,则可通过积分得到新的函数展开式.
	\item 
	若已知某一函数的 Taylor 级数,且可以逐项求导,则对其求导即可得到导函数的 Taylor 级数.
\end{enumerate}
\begin{example}
	$$e^x = \sum_{n=0}^{\infty}\frac{x^n}{n!}$$
\end{example}
\begin{proof}
	$ e^x$的Taylor公式
	\[e^x=1+x+\frac{1}{2!}x^2+\cdots+\frac{1}{n!}x^n+r_n(x).x\in(-\infty,+\infty)\]
	取$r_n(x)$的Lagrange余项
	\[r_n(x)=\frac{e^\xi}{(n+1)!}x^{n+1},\xi\text{介于}x\text{与}0\text{之间}\]
	由\ref{阶的比较}
	\[|r_n(x)|\leqslant \frac{e^{|x|}}{(n+1)!}|x|^{n+1}\to0(n\to\infty)\]
	对一切$x\in(-\infty,+\infty)$成立，从而$$e^x = \sum_{n=0}^{\infty}\frac{x^n}{n!}$$
\end{proof}
常见初等函数如 $\sin x,\cos x,\dfrac{1}{1-x},\dfrac{1}{1+x},\sqrt{1+x},\sqrt{1-x},\ln(1+x),\arctan x$ 等,可先由 Taylor 公式得到展开式,再证明其余项趋于 $0$.


\begin{example}
	由
	\[
	e^x=\sum_{n=0}^{\infty}\frac{x^n}{n!}
	\]
	作代换 $x\mapsto -x$,得
	\[
	e^{-x}=\sum_{n=0}^{\infty}\frac{(-x)^n}{n!}
	=1-x+\frac{x^2}{2!}-\frac{x^3}{3!}+\cdots .
	\]
	因此
	\[
	\frac{e^x+e^{-x}}{2}=\cosh x
	=1+\frac{x^2}{2!}+\frac{x^4}{4!}+\cdots
	=\sum_{n=0}^{\infty}\frac{x^{2n}}{(2n)!},x\in(-\infty,+\infty)
	\]
	同理
	\[
	\sinh x=\frac{e^x-e^{-x}}{2}
	=\sum_{n=0}^{\infty}\frac{x^{2n+1}}{(2n+1)!}.x\in(-\infty,+\infty)
	\]
\end{example}

\begin{example}
	当 $|x|<1$ 时,
	\[
	\frac{1}{1-x}=1+x+x^2+\cdots+x^n+\cdots .
	\]
	令 $x\mapsto -x$,得
	\[
	\frac{1}{1+x}=1-x+x^2-x^3+\cdots+(-1)^n x^n+\cdots
	=\sum_{n=0}^{\infty}(-1)^n x^n.
	\]
	令 $x\mapsto x^2$ 与 $x\mapsto -x^2$,分别有
	\[
	\frac{1}{1-x^2}=1+x^2+x^4+\cdots+x^{2n}+\cdots,
	\]
	\[
	\frac{1}{1+x^2}=1-x^2+x^4-\cdots+(-1)^n x^{2n}+\cdots .
	\]
\end{example}


\begin{example}
	由
	\[
	\frac{1}{1+t}=\sum_{n=0}^{\infty}(-1)^n t^n,
	\qquad |t|<1,
	\]
	逐项积分得
	\[\ln(1+x)
	=\int_{0}^{x}\frac{\dd t}{1+t}\=\int_{0}^{x}\sum_{n=0}^{\infty}(-1)^n t^n\dd t=\sum_{n=0}^{\infty}(-1)^n\frac{x^{n+1}}{n+1}           =x-\frac{x^2}{2}+\frac{x^3}{3}-\cdots+(-1)^n\frac{x^{n+1}}{n+1}+\cdots .
	\]
	收敛域为$(-1,1]$
\end{example}

\begin{example}
	由
	\[
	\frac{1}{1+t^2}=\sum_{n=0}^{\infty}(-1)^n t^{2n},
	\qquad |t|<1,
	\]
	逐项积分得
	\[
	\begin{aligned}
		\arctan x
		&=\int_{0}^{x}\frac{\dd t}{1+t^2}                                      \\
		&=\int_{0}^{x}\sum_{n=0}^{\infty}(-1)^n t^{2n}\dd t                    \\
		&=\sum_{n=0}^{\infty}(-1)^n\frac{x^{2n+1}}{2n+1}                       \\
		&=x-\frac{x^3}{3}+\frac{x^5}{5}-\cdots+(-1)^n\frac{x^{2n+1}}{2n+1}+\cdots .
	\end{aligned}
	\]
	令 $x\to1^{-}$,得
	\[
	\frac{\pi}{4}=1-\frac{1}{3}+\frac{1}{5}-\frac{1}{7}+\cdots,
	\]
	即
	\[
	\pi=4\sum_{n=0}^{\infty}\frac{(-1)^n}{2n+1}.
	\]
	若取前 $n+1$ 项,由 Leibniz 余项估计,
	\[
	|R_n|\leqslant\frac{4}{2n+3}.
	\]
	令 $x=\dfrac{\sqrt{3}}{3}$,由于 $\arctan\dfrac{\sqrt{3}}{3}=\dfrac{\pi}{6}$,有
	\[
	\frac{\pi}{6}
	=\frac{\sqrt{3}}{3}-\frac{1}{3}\left(\frac{\sqrt{3}}{3}\right)^3
	+\cdots+\frac{(-1)^n}{2n+1}\left(\frac{\sqrt{3}}{3}\right)^{2n+1}+\cdots .
	\]
	于是
	\[
	\pi=6\sum_{n=0}^{\infty}\frac{(-1)^n}{2n+1}\left(\frac{\sqrt{3}}{3}\right)^{2n+1}.
	\]
	同样由 Leibniz 余项估计,若取前 $n+1$ 项,则
	\[
	|R_n|\leqslant \frac{6}{2n+3}\left(\frac{\sqrt{3}}{3}\right)^{2n+3}.
	\]
\end{example}


\begin{example}
	由
	\[
	\sin x=\sum_{n=0}^{\infty}(-1)^n\frac{x^{2n+1}}{(2n+1)!},
	\]
	逐项求导可得
	\[
	\cos x=\sum_{n=0}^{\infty}(-1)^n\frac{x^{2n}}{(2n)!}.
	\]
\end{example}

记录几个非零点处的展开.核心方法是把展开点平移到 $0$ 处.

\begin{example}
	对 $\sin x$ 在 $x=\dfrac{\pi}{2}$ 处展开.因为
	\[
	\sin x=\sin\left(x-\frac{\pi}{2}+\frac{\pi}{2}\right)=\cos\left(x-\frac{\pi}{2}\right),
	\]
	故
	\[
	\sin x=\sum_{n=0}^{\infty}(-1)^n\frac{\left(x-\frac{\pi}{2}\right)^{2n}}{(2n)!}.
	\]
\end{example}

\begin{example}
	在 $x=\dfrac{\pi}{4}$ 处,
	\[
	\begin{aligned}
		\sin x
		&=\sin\left(x-\frac{\pi}{4}+\frac{\pi}{4}\right)                      \\
		&=\sin\left(x-\frac{\pi}{4}\right)\cos\frac{\pi}{4}
		+\cos\left(x-\frac{\pi}{4}\right)\sin\frac{\pi}{4}                     \\
		&=\frac{\sqrt{2}}{2}\sum_{n=0}^{\infty}(-1)^n\frac{\left(x-\frac{\pi}{4}\right)^{2n+1}}{(2n+1)!}
		+\frac{\sqrt{2}}{2}\sum_{n=0}^{\infty}(-1)^n\frac{\left(x-\frac{\pi}{4}\right)^{2n}}{(2n)!}.
	\end{aligned}
	\]
\end{example}
设 $f(x)$ 的幂级数展开为
\[
f(x)=\sum_{n=0}^{\infty}a_nx^n,
\]
收敛半径为 $R_1$;$g(x)$ 的幂级数展开为
\[
g(x)=\sum_{n=0}^{\infty}b_nx^n,
\]
收敛半径为 $R_2$.则 $f(x)g(x)$ 的幂级数展开就是它们的 Cauchy 乘积:
\[
\begin{aligned}
	f(x)g(x)
	&=\left(\sum_{n=0}^{\infty}a_nx^n\right)
	\left(\sum_{n=0}^{\infty}b_nx^n\right)        \\
	&=\sum_{n=0}^{\infty}c_nx^n,
\end{aligned}
\]
其中
\[
c_n=\sum_{k=0}^{n}a_kb_{n-k},
\qquad |x|<\min\{R_1,R_2\}.
\]

当 $b_0\neq0$ 时,可以通过待定系数法求 $\dfrac{f(x)}{g(x)}$ 的幂级数展开.设
\[
\frac{f(x)}{g(x)}=\sum_{n=0}^{\infty}c_nx^n,
\]
则
\[
\left(\sum_{n=0}^{\infty}b_nx^n\right)
\left(\sum_{n=0}^{\infty}c_nx^n\right)
=\sum_{n=0}^{\infty}a_nx^n.
\]
分离 $x$ 的各次幂的系数,可依次得到
\[
\begin{alignedat}{5}
	b_0c_0&=a_0
	&\qquad&\Longrightarrow
	&\qquad c_0&=\frac{a_0}{b_0},\\
	b_0c_1+b_1c_0&=a_1
	&\qquad&\Longrightarrow
	&\qquad c_1&=\frac{a_1-b_1c_0}{b_0},\\
	b_0c_2+b_1c_1+b_2c_0&=a_2
	&\qquad&\Longrightarrow
	&\qquad c_2&=\frac{a_2-b_1c_1-b_2c_0}{b_0},\\
	&\cdots
\end{alignedat}
\]
一直继续下去,可求得所有的 $c_n$.

\begin{example}[例 10.4.4]
	求 $e^x\sin x$ 在 $x=0$ 的幂级数展开,到 $x^5$.
\end{example}

\begin{solution}
	由 $e^x$ 与 $\sin x$ 在 $x=0$ 处的展开式,有
	\[
	\begin{aligned}
		e^x\sin x
		&=\left(1+x+\frac{x^2}{2!}+\frac{x^3}{3!}+\frac{x^4}{4!}+\cdots\right)
		\left(x-\frac{x^3}{3!}+\frac{x^5}{5!}-\cdots\right)\\
		&=x+x^2+\frac{1}{3}x^3-\frac{1}{30}x^5+\cdots .
	\end{aligned}
	\]
	其中 $x^4$ 项的系数为 $0$.上述幂级数展开对一切 $x\in(-\infty,+\infty)$ 都成立.
\end{solution}

\begin{example}
	求 $\tan x$ 在 $x=0$ 的幂级数展开,到 $x^5$.
\end{example}

\begin{solution}
	由于 $\tan x$ 是奇函数,我们可以令
	\[
	\tan x=\frac{\sin x}{\cos x}=c_1x+c_3x^3+c_5x^5+\cdots .
	\]
	于是
	\[
	\left(c_1x+c_3x^3+c_5x^5+\cdots\right)
	\left(1-\frac{x^2}{2!}+\frac{x^4}{4!}-\cdots\right)
	=x-\frac{x^3}{3!}+\frac{x^5}{5!}-\cdots .
	\]
	比较等式两端 $x,x^3$ 与 $x^5$ 的系数,可得
	\[
	\begin{cases}
		c_1=1,\\[4pt]
		c_3-\dfrac{c_1}{2!}=-\dfrac{1}{3!},\\[8pt]
		c_5-\dfrac{c_3}{2!}+\dfrac{c_1}{4!}=\dfrac{1}{5!}.
	\end{cases}
	\]
	因而
	\[
	c_1=1,
	\qquad
	c_3=\frac{1}{3},
	\qquad
	c_5=\frac{2}{15}.
	\]
	因此
	\[
	\tan x=x+\frac{1}{3}x^3+\frac{2}{15}x^5+\cdots .
	\]
\end{solution}

\begin{example}
	求 $\ln \dfrac{\sin x}{x}$ 在 $x=0$ 的幂级数展开,到 $x^4$,其中函数 $\dfrac{\sin x}{x}$ 应理解为
	\[
	f(x)=
	\begin{cases}
		\dfrac{\sin x}{x},&x\neq0,\\[6pt]
		1,&x=0.
	\end{cases}
	\]
\end{example}

\begin{solution}
	首先,利用 $\sin x$ 的幂级数展开,可以得到
	\[
	\frac{\sin x}{x}=1-\frac{x^2}{3!}+\frac{x^4}{5!}-\cdots .
	\]
	另外,我们有
	\[
	\ln(1+u)=u-\frac{u^2}{2}+\frac{u^3}{3}-\cdots .
	\]
	将
	\[
	u=-\frac{x^2}{3!}+\frac{x^4}{5!}-\cdots
	\]
	代入上式,即得
	\[
	\begin{aligned}
		\ln\frac{\sin x}{x}
		&=\left(-\frac{x^2}{3!}+\frac{x^4}{5!}-\cdots\right)
		-\frac{1}{2}\left(-\frac{x^2}{3!}+\frac{x^4}{5!}-\cdots\right)^2+\\
		&\quad+\frac{1}{3}\left(-\frac{x^2}{3!}+\frac{x^4}{5!}-\cdots\right)^3-\cdots\\
		&=-\frac{x^2}{6}+\left(\frac{1}{120}-\frac{1}{72}\right)x^4+\cdots\\
		&=-\frac{x^2}{6}-\frac{x^4}{180}+\cdots .
	\end{aligned}
	\]
	
且
	\[
	\frac{\sin x}{x}=\prod_{n=1}^{\infty}\left(1-\frac{x^2}{n^2\pi^2}\right).
	\]
	两边取对数,再将 $\ln\left(1-\dfrac{x^2}{n^2\pi^2}\right)$ 展开成幂级数,有
	\[
	\begin{aligned}
		\ln\frac{\sin x}{x}
		&=\sum_{n=1}^{\infty}\ln\left(1-\frac{x^2}{n^2\pi^2}\right)\\
		&=-\sum_{n=1}^{\infty}\left(\frac{x^2}{n^2\pi^2}+\frac{1}{2}\frac{x^4}{n^4\pi^4}+\cdots\right).
	\end{aligned}
	\]
	比较它们的 $x^2$ 系数与 $x^4$ 系数应对应相等,于是得到
	\[
	\sum_{n=1}^{\infty}\frac{1}{n^2}=\frac{\pi^2}{6},
	\qquad
	\sum_{n=1}^{\infty}\frac{1}{n^4}=\frac{\pi^4}{90}.
	\]
	如果我们计算时更加细致,也就是将 $\ln\dfrac{\sin x}{x}$ 的幂级数展开计算到 $x^6,x^8,\cdots$,还可以获得
	\[
	\sum_{n=1}^{\infty}\frac{1}{n^6},
	\quad
	\sum_{n=1}^{\infty}\frac{1}{n^8},
	\quad\cdots
	\]
	的精确值.
\end{solution}
最后我们给出一个常用Taylor级数的表

在$x=0$的邻域内，有
\begin{enumerate}
	\item \[
	e^x =\sum_{n = 0}^{\infty}\frac{x^n}{n!}= 1 + x + \frac{1}{2!}x^2 + \frac{1}{3!}x^3 + \cdots + \frac{1}{n!}x^n +\cdots 
	,x\in(-\infty,+\infty)\]
	\[
	\frac{e^x+e^{-x}}{2}=\cosh x
	=1+\frac{x^2}{2!}+\frac{x^4}{4!}+\cdots
	=\sum_{n=0}^{\infty}\frac{x^{2n}}{(2n)!},x\in(-\infty,+\infty)
	\]
	\[
	\sinh x=\frac{e^x-e^{-x}}{2}
	=\sum_{n=0}^{\infty}\frac{x^{2n+1}}{(2n+1)!}.x\in(-\infty,+\infty)
	\]
	\item \[
	\sin x=\sum_{n=0}^{\infty}(-1)^n\frac{x^{2n+1}}{(2n+1)!},x\in(-\infty,+\infty)
	\]
	\[
	\cos x=\sum_{n=0}^{\infty}(-1)^n\frac{x^{2n}}{(2n)!}.x\in(-\infty,+\infty)
	\]
	\[
	\tan x = x + \frac{1}{3}x^3 + \frac{2}{15}x^5 + \frac{17}{315}x^7 + \cdots
	\]
	\item 
	若$ \alpha \notin \mathbb{Z}_+$, $\forall k \in \mathbb{R}$, $C_\alpha^k \neq 0$,
	\[(1+x)^\alpha=\sum_{n=0}^{\infty}\binom{\alpha}{n} x^n, \quad \begin{cases}x \in(-1,1), & \text { 当 } \alpha \leqslant-1, \\ x \in(-1,1], & \text { 当 }-1<\alpha<0, \\ x \in[-1,1], & \text { 当 } \alpha>0 .\end{cases}\]	
	其中\[\binom{\alpha}{n}=\frac{\alpha(\alpha-1)\cdots(\alpha-n+1)}{n!}\]
	
	当 $\alpha = -1$ 时，
	\[
	\frac{1}{1+x} =\sum_{x=0}^{\infty}(-1)^nx^n= 1 - x + x^2 - x^3 + x^4 - \cdots + (-1)^n x^n +\cdots
	,x\in(-1,1)\]
	
	对于 $\alpha = \dfrac{1}{2}$，
	\[\begin{aligned}
		\sqrt{1+x} &=1 + \frac{1}{2}x+\sum_{k=2}^{\infty}\dfrac{(2k-3)!!}{(2k)!!}\\
		&= 1 + \frac{1}{2}x - \frac{1}{2 \cdot 4}x^2 + \frac{1 \cdot 3}{2 \cdot 4 \cdot 6}x^3 - \cdots + (-1)^{n-1} \frac{(2n-3)!!}{(2n)!!}x^n + \cdots
	\end{aligned}
	\]
	对于 $\alpha = -\dfrac{1}{2}$，
	\[\begin{aligned}
		\frac{1}{\sqrt{1+x}}&=\sum_{n=0}^{\infty}(-1)^n \frac{(2n-1)!!}{(2n)!!}\\&= 1 - \frac{1}{2}x + \frac{1 \cdot 3}{2 \cdot 4}x^2 - \frac{1 \cdot 3 \cdot 5}{2 \cdot 4 \cdot 6}x^3 + \cdots + (-1)^n \frac{(2n-1)!!}{(2n)!!}x^n +\cdots
	\end{aligned}
	\]
	使用代换法
	\[
	\frac{1}{1-x} =\sum_{x=0}^{\infty}x^n 
	\]
	\[
	\frac{1}{1+x^2} = \sum_{x=0}^{\infty}(-1)^nx^{2n}
	\]
	\[
	\frac{1}{1-x^2} = \sum_{x=0}^{\infty}x^{2n}
	\]
	\item 
	\[
	\ln(1+x) = \sum_{n=0}^{\infty}(-1)^n\frac{x^{n+1}}{n+1}=x - \frac{1}{2}x^2 + \frac{1}{3}x^3 - \cdots + (-1)^{n-1} \frac{1}{n} x^n + \cdots,x\in(-1,1]
	\]
	\[
	\ln(1-x) =\sum_{n=0}^{\infty}(-1)\frac{x^{n+1}}{n+1}= -x - \frac{1}{2}x^2 - \frac{1}{3}x^3 - \cdots - \frac{1}{n}x^n + \cdots
	x\in[-1,1)\]
	\[
	\tan^{-1} x=\sum_{n=1}^{\infty}\frac{x^{2n+1}}{(2n+1)!} = x - \frac{1}{3}x^3 + \frac{1}{5}x^5 - \frac{1}{7}x^7 + \cdots + (-1)^n \frac{1}{2n+1} x^{2n+1} + \cdots,x\in[-1,1]
	\]
	\[\begin{aligned}
		\sin^{-1} x&=x+\sum_{n=1}^{\infty}\frac{(2n-1)!!}{(2n)!!}\frac{x^{2n+1}}{2n+1}\\ &= x + \frac{1}{2} \cdot \frac{1}{3} x^3 + \frac{1 \cdot 3}{2 \cdot 4} \cdot \frac{1}{5} x^5 + \cdots + \frac{(2n-1)!!}{(2n)!!} \cdot \frac{1}{2n+1} x^{2n+1} + \cdots,x\in[-1,1]
	\end{aligned}
	\]
	而
	\[
	\cos^{-1} x = \frac{\pi}{2} - \sin^{-1} x
	\]
\end{enumerate}
